%! Author = Эчпочмак
%! Date = 09.05.2023

% Packages
\usepackage{amsmath}
\usepackage[utf8]{inputenc}
\usepackage[T1]{fontenc}

% Document
\begin{document}

Извините, пожалуйста, похоже, что в моем предыдущем ответе произошла ошибка. Верный ответ на задачу sqrt(x^2-2x+1)<=sqrt(x^2+x) можно получить следующим образом:

\begin{align*} \sqrt{x^2-2x+1} &\leq \sqrt{x^2+x} \ |x-1| &\leq \sqrt{x(x+1)} \end{align*}

Дальше нужно рассмотреть два случая в зависимости от значения $x$.

Пусть $x \geq 0$. Тогда: \[x-1 \leq \sqrt{x(x+1)} \text{ или } -x+1 \leq \sqrt{x(x+1)}\]

Решая первое неравенство, получаем: \[ \begin{aligned} x-1 &\leq \sqrt{x(x+1)} \ x^2-2x+1 &\leq x^2+x \ 2x &\geq -1 \ x &\geq -\frac{1}{2} \end{aligned} \]

Решая второе неравенство, получаем: \[ \begin{aligned} -x+1 &\leq \sqrt{x(x+1)} \ x^2-2x+1 &\geq x^2+x \ -3x &\geq -1 \ x &\leq \frac{1}{3} \end{aligned} \]

Таким образом, если $x \geq 0$, то решением неравенства является множество $[-\frac{1}{2}, \frac{1}{3}]$.

Пусть $x < 0$. В этом случае $|x-1|=1-x$, поэтому изначальное неравенство примет вид: \[ 1-x \leq \sqrt{x(x+1)} \]

Возводим все в квадрат и приводим подобные слагаемые: $$ x^2-2x+



\end{document}